\section{A transverse saving in every fixed dimension}
\label{t:section}

Throughout this section the spatial dimension $n\ge3$ is fixed, and
\[
 \beta=\frac{2}{n+1},\qquad e=\frac{n^2}{n+1}.
\]
For a frequency $\xi$ write $v(\xi)=(2\xi,1)$; replacing these
vectors by unit vectors changes determinants by bounded factors on
a fixed bounded frequency region.  We write $Ef=\mathcal E f$ for the
exact extension operator of Section~\ref{d:section}.  The following statement concerns
exact extension fields.  Its use for smooth packet arrays will come
after the exact Fourier-layer reduction in the next section.

\begin{proposition}[Transverse saving]\label{t:transverse}
Fix a bounded frequency region in $\R^n$.  There are $c>0$ and
$L_0<\infty$, depending only on this region and $n$, such that the
following configuration is impossible for $L\ge L_0$.  Let
$\mathcal Q$ be distinct unit lattice cubes in a ball of radius $L^2$
in $\R^{n+1}$, with
\begin{equation}\label{t:forbidden-counts}
 \#\mathcal Q\ge L^{2n-c},\qquad
 \#\{q\in\mathcal Q:q\cap B(z,r)\ne\varnothing\}
       \le L^c r^n\quad(r\ge1).
\end{equation}
Suppose $f_1,\ldots,f_{n+1}$ have $L^2$ norm at most one, frequency
support in the fixed region, and
\[
 \sup_q|Ef_i|\ge L^{-e-c}
       \quad(q\in\mathcal Q,\ 1\le i\le n+1).
\]
It is impossible that every choice of one frequency from each support
has direction determinant at least $L^{-c}$ in absolute value.
\end{proposition}

We prove the equivalent assertion that no countersequence exists with
$c\downarrow0$ and $L\to\infty$.  In this section $o(1)$ refers to
such a sequence, after subsequences when necessary.  Estimates available
with every fixed positive power loss may be diagonalized to have
$L^{o(1)}$ loss.  The constants in each such estimate are fixed before
diagonalization.  The proof uses the limiting information calculus of
\citet[Section~3, ``Limiting information calculus'']{PSE2026}.  We state
the usable facts below and prove the changes needed in dimension $n$.
In particular, the planar transverse proposition in that source is not
used as a higher-dimensional assertion.

\subsection{Uniform multilinear estimates and incident weights}

We first specify the dependence on transversality.  This prevents a
subpower determinant loss from concealing an uncontrolled constant.

\begin{lemma}[Multilinear estimates with polynomial dependence]
\label{t:multilinear}
If $n+1$ bounded-frequency exact inputs are transverse with determinant
at least $\nu\in(0,1]$, then, for each $\eps>0$,
\begin{equation}\label{t:multilinear-extension}
 \sum_{q\subset B_S}\prod_{i=1}^{n+1}
       \bigl(\sup_q|Ef_i|\bigr)^{2/n}
 \le C_\eps\nu^{-M_\eps}S^\eps
                    \prod_{i=1}^{n+1}\norm{f_i}_2^{2/n}.
\end{equation}
The same bound holds with the sum replaced by the integral of the
product.  For transverse families of tubes of width $\delta$ in a fixed
bounded region, with nonnegative weights $w_i(T)$, the corresponding
Kakeya bound is
\begin{equation}\label{t:kakeya}
 \int\prod_{i=1}^{n+1}
       \left(\sum_Tw_i(T)\mathbf1_T\right)^{1/n}
 \le C_\eps\nu^{-M_\eps}\delta^{n+1-\eps}
             \prod_i\left(\sum_Tw_i(T)\right)^{1/n}.
\end{equation}
Fixed enlargements are allowed.  Subpower enlargements and
$\nu=S^{-o(1)}$ give only subpower losses.
\end{lemma}

\begin{proof}
Subdivide direction supports into neighborhoods of radius a sufficiently
small fixed power of $\nu$.  There are $\nu^{-O_n(1)}$ choices of
$(n+1)$ neighborhoods.  For each choice, send representative directions
to the coordinate axes.  The map, its inverse, and its determinant have
size bounded by fixed powers of $\nu^{-1}$.  The near-coordinate-axes
multilinear Kakeya theorem
\citep[Theorem~1.15 and Remark~1.10]{BCT2006}, or its strengthened form
\citep{Guth2010}, gives \eqref{t:kakeya} after summing the choices.
Weights follow by repetition and approximation.

Here is the finite induction giving the required uniform restriction
constant.  Let $A_B(S,\nu)$ be the best normalized integral constant
for inputs in a fixed box of size $B$.  Decompose each input into exact
packets at radius $S$, and on each $\sqrt S$-box retain the packets
meeting its $S^\theta\sqrt S$ enlargement.  Their input squared norms
are bounded by incident coefficient energy, and the discarded fields
are $O_N(S^{-N})$ relative to their parent input norms.  The enlarged
supports remain $\nu/2$ transverse if $S^{-1/2}\le c_B\nu$.
After rescaling by $S$, the padded tubes have width
$\delta=S^{-1/2+\theta}$.  Apply \eqref{t:kakeya} with loss $\sigma$.
Dividing its bound by the volume $S^{-(n+1)/2}$ of a rescaled
$\sqrt S$-box gives
\[
 A_B(S,\nu)\le C\nu^{-M_\sigma}S^a
          A_{B+1}(C_B\sqrt S,\nu/2)+O_N(S^{-N}),
 \qquad a=(n+1)\theta+\sigma/2.
\]
Increasing $N$ absorbs all errors into the same inequality for $1+A_B$;
when expanding products with exponent $2/n\le1$, use subadditivity to
bound terms containing a discarded factor.  For a prescribed $\eps$,
choose a finite $J$ with $(n+1)2^{-J}<\eps/2$, and then $\theta,\sigma$
with $2a<\eps/2$.  Iterating $J$ times and using the terminal trivial
bound $O(S_J^{n+1})$, where $S_J\asymp S^{2^{-J}}$, proves the
integral estimate when $S\ge(C/\nu)^{2^J}$.  In the complementary
range the trivial bound is absorbed by increasing $M_\eps$.
Thus the dependence on $\nu^{-1}$ is a fixed power.

For the cube-supremum version the exponent $p=2/n$ can be below one,
so we include the submean argument.  If $F$ has Fourier support in the
fixed compact region, choose a Schwartz reproducing kernel $\psi$ and
write $F=\psi*F$.  For $A>0$ let
$M_AF(x)=\sup_y|F(x-y)|(1+|y|)^{-A}$, which is finite because these
extensions are bounded.  In the convolution insert
$|F(w)|\le|F(w)|^p M_AF(x)^{1-p}(1+|x-w|)^{A(1-p)}$.
Peetre's elementary inequality
$(1+|x-w|)\le(1+|x-z|)(1+|z-w|)$ then gives, uniformly in $z$,
\[
 \frac{|F(z)|}{(1+|x-z|)^A}
 \le C_A M_AF(x)^{1-p}
             \int(1+|x-w|)^{-Ap}|F(w)|^p\dd w,
\]
using the boundedness of $(1+|z-w|)^A|\psi(z-w)|$.
Taking the supremum and dividing by $M_AF(x)^{1-p}$ proves the
required Schwartz-average bound for $M_AF(x)^p$; the zero case is
immediate.  Choose $Ap$ arbitrarily large.  In particular this bounds
$(\sup_q|Ef_i|)^{2/n}$ by that average for each $x\in q$.
Integrate the product over $q$ and sum.  For each independent list of
shifts the translated fields are exact extensions with unchanged norms
and supports.  Apply the integral estimate and integrate the Schwartz
weights.  This proves \eqref{t:multilinear-extension}.
\end{proof}

\begin{lemma}[Pruned incident-tube model]\label{t:tubes}
A countersequence to Proposition~\ref{t:transverse} yields a set
$\mathcal P$ of centers of distinct $L^{-1}$-scale cubes in a bounded
region, satisfying
\begin{equation}\label{t:point-counts}
 \#\mathcal P=L^{n+o(1)},\qquad
 \#(\mathcal P\cap B(z,r))\le L^{n+o(1)}r^n
                     \quad(L^{-1}\le r\le1),
\end{equation}
and $n+1$ transverse weighted tube families of width $L^{-1+o(1)}$.
Their weights satisfy $\sum_Tw_i(T)\lesssim1$.  For every $x\in\mathcal P$
there is an incident subfamily of retained weight
$w_i^*(x)=L^{-e+o(1)}$ whose normalized law satisfies
\begin{equation}\label{t:strip-law}
 \mathbb P\{\xi(T_i)\in N_{L^{-\alpha}}(H)\mid X=x\}
       \le L^{-c_1\alpha+o(1)}
\end{equation}
for every affine $(n-2)$-plane $H\subset\R^n$ and every fixed rational
$0<\alpha<1/8$.  Here $c_1>0$ is fixed independently of $\alpha$, and
the estimate is uniform in $x,H$ for each fixed $\alpha$.
\end{lemma}

\begin{proof}
Sort the original unit cubes by their occupancy in $L$-boxes.  A dyadic
class retains $L^{2n-o(1)}$ cubes, on which the squared supremum sum of
each field is at least $L^{2n-2e-o(1)}=L^{2n/(n+1)-o(1)}$.
The refined fractal estimate, Proposition~\ref{d:refined}, at radius
$L^2$ improves this upper bound by a fixed power if the $L$-box
occupancy is at most $L^{n-\delta}$ for any fixed $\delta>0$.
The counting hypothesis supplies the upper bound $L^{n+o(1)}$.
Consequently the retained occupancies are $L^{n+o(1)}$.
Divide coordinates by $L^2$ and use the occupied box centers.  Counting
the original cubes in an enlarged ball, then dividing by the minimum
occupancy, proves \eqref{t:point-counts}.

Exact packet localization at radius $L^2$ gives coefficient weights with
bounded total sum.  On each retained physical $L$-box, keep trajectories
in a padded enlargement.  The padding exponent can tend sufficiently
slowly to zero; rapid tails make the discarded field negligible.
The retained exact input has squared norm at most $Cw_i(x)$, where
$w_i(x)$ is its incident weight.  Apply Proposition~\ref{d:fractal}
at radius $L$ to its $L^{n+o(1)}$ selected unit cubes:
\[
 L^{n-2e-o(1)}\le L^{n/(n+1)+o(1)}w_i(x),
 \qquad w_i(x)\ge L^{-e-o(1)}.
\]
The counting version of \eqref{t:kakeya} gives
$\sum_x\prod_iw_i(x)^{1/n}\le L^{o(1)}$.  If some weight exceeds
$L^{-e+\delta}$, this and the other lower bounds allow only
$L^{n-\delta/n+o(1)}$ such points.  Diagonal deletion leaves
$w_i(x)=L^{-e+o(1)}$ uniformly.

Let $c_0$ be the constant in the codimension-two strip gain,
Proposition~\ref{d:strip}, and fix $0<c_1<c_0/2$.  For a finite list of
rational $\alpha$, successively remove, at each $x,i$, strips of width
$L^{-\alpha}$ carrying remaining weight greater than
$L^{-c_1\alpha}w_i(x)$.  At one scale there are at most
$L^{c_1\alpha}$ removals.  Assign a packet to its first removal, so
the corresponding subarray energies sum to $O(w_i(x))$.  On the
physical $L$-box the exact support enlargement is $O(L^{-1})$, smaller
than the strip width.  Proposition~\ref{d:strip}, followed by Cauchy
over the removed strips, bounds their combined squared supremum sum by
\[
 L^{n/(n+1)-(c_0-c_1)\alpha+o(1)}w_i(x).
\]
This is negligible compared with the required local squared norm.
The triangle inequality in that square-sum norm, followed by the
unrestricted local fractal estimate, forces surviving weight
$L^{-e+o(1)}$.  On termination every affine strip has remaining weight
at most the deletion threshold.  Normalizing by the surviving weight
costs $L^{o(1)}$.  Diagonalize over finite rational lists to obtain
\eqref{t:strip-law}.  Constant or subpower enlargements of the strips
are covered by applying the assertion at a slightly smaller rational
exponent.  The pruning is allowed to depend on $x$; it defines the
conditional laws used next.
\end{proof}

\subsection{Information profiles and the cost of direction bins}

Draw $X$ uniformly from $\mathcal P$.  Conditional on $X$, draw
$T_i,T_i'$ independently from the retained law of family $i$, for
$1\le i\le n+1$.  A tube label contains its central line and exact
direction.  Let $X_s$ be its position label in nested cube grids of
side comparable to $L^{-s}$, for rational $s\in[0,1]$, with the
trivial partition at zero.  For a discrete variable $Y$ and conditioning
$V$, information at the sampled labels means
\[
 I_L(Y\mid V)=-\log_L\mathbb P(Y=Y(\omega)\mid V).
\]
All spatial, coordinate, and direction partitions below have at most
$L^{C+o(1)}$ labels.  Exact tube labels may be more numerous; only
their role as conditioning variables is used.

\begin{lemma}[Limiting information calculus]\label{t:information}
After passing to a subsequence, the countably many information variables
at rational block and prefix scales have a common limiting joint law.
Their conditional informations are uniformly integrable and bounded in
the limit.  The chain rule is exact, and conditioning lowers information
almost surely in this law.  If, on a specified event and given $(Z,V)$,
there are at most $L^{a+o(1)}$ possible labels for $Y$, then the limiting
information $I(Y\mid Z,V)$ is at most $a$ on that event.  Event
indicators can be included in the joint law; finite systems of limiting
inequalities of positive probability pass to finite scale with fixed
positive slack.
\end{lemma}

\begin{proof}
This is the dimension-free lemma with this name in
\citet[Section~3]{PSE2026}.  The short argument also explains its
pointwise scope.  If a variable has at most $L^{C+o(1)}$ labels, then
$\mathbb P\{I_L(Y\mid V)>C+a\}\le L^{-a+o(1)}$, uniformly in
conditioning; integrating this bound gives uniform integrability.
Writing $p=\mathbb P(Y\mid V)$ and $q=\mathbb P(Y\mid V,W)$ at
the sampled labels, the conditional expectation of $p/q$ is at most
one.  Thus
\[
 \mathbb P\{I_L(Y\mid V)-I_L(Y\mid V,W)<-a\}\le L^{-a}.
\]
This proves the limiting conditioning inequality, without asserting
finite-scale pointwise monotonicity.  The restricted determination rule
follows by summing probabilities over the permitted small label set.
Tightness and a diagonal subsequence give the joint law.  Include all
rational tests, their costs, and the restricted-event indicators;
open-set implications with slack justify their passage to the limit.
\end{proof}

The information profiles are
\[
 F(s)=I(X_s),\qquad G_i(s)=I(X_s\mid T_i),\qquad
 G_i'(s)=I(X_s\mid T_i').
\]
Thus $F$ records the limiting exponent of the reciprocal probability
of the sampled spatial cell, whereas $G_i$ records that exponent after
one incident tube is specified.  We shall compare these profiles in
coordinates adapted to the tube directions; their endpoint values
will force equality in that comparison.

These profiles extend to nondecreasing Lipschitz functions on $[0,1]$,
with constants $n+1$ and $1$, respectively, and vanish at zero.  Indeed
a tube intersects at most $L^{t+o(1)}$ children inside a parent when
the exponent increases by $t$.  The determination rule gives the
conditional Lipschitz constant; unrestricted grid counts give the
other one.  Moreover
\begin{equation}\label{t:profile-endpoints}
 F(s)\ge ns,\qquad F(1)=n,\qquad
 G_i(1),G_i'(1)\ge\frac n{n+1}.
\end{equation}
The first two assertions follow from \eqref{t:point-counts}.  For the
third, the retained weights give
\[
 \mathbb P(X=x,T_i=T)\le L^{-n+e+o(1)}w_i(T)
       =L^{-n/(n+1)+o(1)}w_i(T).
\]
The probability of sampling a tube with marginal probability smaller
than $L^{-a}w_i(T)$ is at most $L^{-a}\sum_Tw_i(T)=O(L^{-a})$.
Outside this event the conditional point mass is at most
$L^{-n/(n+1)+a+o(1)}$.  Let $a\downarrow0$, also for the primed tubes.

To choose coordinates adapted to reference tubes, we reveal their
coarse direction bins.  Conditioning on these bins may lower the
limiting information; the costs below measure this loss.
Fix a rational block $[s,s+b]$, a rational $b_*\ge b$, and $C=X_s$.
Let $D$ contain direction bins of precision $L^{-b_*}$ for a fixed
sublist of the $2(n+1)$ tubes, including one reference tube per family.
Normalize the parent to size one and use the basis given by deterministic
representatives of the reference bins.  Transversality bounds its norm
and inverse by $L^{o(1)}$.  For a prefix $0<t\le b$, let
$Z=(Z_1,\ldots,Z_{n+1})$ be the coordinate grid labels at resolution
$L^{-t}$.  Adding bins to $D$ does not change this basis or these labels.

\begin{lemma}[Block comparison, additivity, and coordinate subsets]
\label{t:additivity}
Put $J=F(s+t)-F(s)$ and $J_i=G_i(s+t)-G_i(s)$ for the selected references.
The nonnegative costs
\[
 \ell=I(X_{s+t}\mid C)-I(X_{s+t}\mid C,D),\quad
 \ell_i=I(X_{s+t}\mid C,T_i)-I(X_{s+t}\mid C,T_i,D)
\]
satisfy
\begin{equation}\label{t:block-comparison}
 \sum_i(J_i-\ell_i)
 \le\sum_i I(Z_i\mid Z_{\widehat i},C,D)
 \le I(Z\mid C,D)=J-\ell.
\end{equation}
For disjoint blocks or prefixes with the same bin list and the same
exact tube conditioning, the sum of expectations of each cost is
$O_n(b_*)$.  Almost surely, for every choice of references,
\begin{equation}\label{t:replacement}
 F=\sum_{i=1}^{n+1}G_i,\qquad G_i'=G_i,\qquad
 G_i(1)=\frac n{n+1}.
\end{equation}
For every coordinate subset $A\subset\{1,\ldots,n+1\}$,
\begin{equation}\label{t:subset-information}
 I((Z_i)_{i\in A}\mid C,D)
       =\sum_{i\in A}J_i+O_n\left(\sum_i\ell_i\right).
\end{equation}
No unused coordinates are included in the conditioning in this formula.
\end{lemma}

\begin{proof}
A reference tube determines its off-axis coordinates to $L^{o(1)}$
choices: its relative width and its direction-bin error are both at
most $L^{-t+o(1)}$.  Given $(C,D)$, the block child label
$X_{s+t}$ and the coordinate label $Z$ determine each other up to
$L^{o(1)}$ choices: the bin-defined change of basis has inverse norm
$L^{o(1)}$, and grid-boundary ambiguity has only subpower multiplicity.
Thus
\[
 J_i-\ell_i=I(Z_i\mid C,D,T_i)
 =I(Z_i\mid Z_{\widehat i},C,D,T_i)
 \le I(Z_i\mid Z_{\widehat i},C,D).
\]
In any ordering of the coordinates each chain-rule term conditions on
a subset of $Z_{\widehat i}$, proving \eqref{t:block-comparison}.

For a finite-valued label $D$, its Shannon entropy is
$H(D)=\mathbb E[-\log p_D(D)]$, with the natural logarithm.
At finite scale the expected cost is the conditional mutual information
of child and $D$ given the parent, divided by $\log L$, with the same
fixed exact tube additionally conditioned on if present.  These mutual
informations telescope along the spatial filtration.  Filling gaps
between disjoint blocks only adds nonnegative terms.  The total is at
most $H(D)/\log L=O_n(b_*)$, since $D$ has $L^{O_n(b_*)}$ labels.
Uniform integrability passes this to the limiting law.  This is exactly
the hypothesis of the ``Block comparison and telescoping costs'' lemma
in \citet[Section~3]{PSE2026}; the number of reference directions only
changes the dimensional constant.

Partition a rational interval into blocks with $b=b_*\downarrow0$.
The expected positive excess of the summed $G_i$ increments over the
$F$ increment is $O_n(b)$, so in the limit every interval has
$\sum_i\Delta G_i\le\Delta F$.  At $[0,1]$, \eqref{t:profile-endpoints}
forces equality.  Nonnegative additive deficits must therefore vanish
on every subinterval.  Repeat this argument for all finitely many
reference replacements and subtract the identities.  This proves
\eqref{t:replacement}.

Finally choose any ordering that puts $A$ first.  Each chain-rule term
is at least $J_i-\ell_i$, and their total surplus is
$\sum_i\ell_i-\ell\le\sum_i\ell_i$ by additivity.  Summing just the
terms for $A$ proves \eqref{t:subset-information}.  Equivalently, the
lower bound for the complement subtracted from the total gives the
upper bound for $A$.  This is why a coordinate pair can be treated
without fixing the other $n-1$ coordinates.
\end{proof}

\subsection{Codimension two and a single common profile}

The pruning in Lemma~\ref{t:tubes} was against affine $(n-2)$-planes,
rather than just angular balls.  The following linear algebra is the
reason for that choice.

\begin{lemma}[Restriction of two direction equations]\label{t:rank}
Let a basis of $\R^{n+1}$ consist of bounded vectors $v(\xi_i)$ with
inverse norm $L^{o(1)}$.  Normalize the $i$th coefficient of another
bounded vector $v(\xi)$ to one, assuming it and its inverse are
$L^{o(1)}$, and denote two of its other coefficient ratios by $a_j,a_k$.
A ball of radius $L^{-u}$ in these two ratios pulls back, on the bounded
frequency region, into an $L^{-u+o(1)}$-neighborhood of an affine
$(n-2)$-plane.  The same conclusion holds when both specified ratios
are $L^{-u}$-small.
\end{lemma}

\begin{proof}
Let $r_\ell$ be the dual basis rows.  Equations for ratios centered at
$(a_j,a_k)$ are the homogeneous equations
$(r_j-a_jr_i)\cdot v=0$ and $(r_k-a_kr_i)\cdot v=0$.
Their row map has smallest singular value $L^{-o(1)}$: the coefficients
of $r_j,r_k$ in any linear combination are the coefficients of that
combination itself, and the primal and dual basis norms are subpower.
Choose the center at a realized ratio in the ball, enlarging its radius
by a fixed factor.  Both rows then vanish at a bounded vector
$v_0=(2\xi_0,1)$.  If a linear combination is $(a,b)$, this implies
$b=-2a\cdot\xi_0$, so $\norm{(a,b)}\lesssim\norm a$.
Consequently the spatial restrictions of the two rows still have
smallest singular value $L^{-o(1)}$.  Their time-one slice is an affine
$(n-2)$-plane, and approximate equations give distance
$L^{-u+o(1)}$ from it.  For two small ratios one can center the equations
at zero; they vanish at the $i$th reference direction, giving the same
argument.  This verifies both rank and its quantitative conditioning
after restriction to the time-one hyperplane.
\end{proof}

\begin{lemma}[One common profile]\label{t:common}
Almost surely all $n+1$ conditional profiles coincide:
\begin{equation}\label{t:common-profile}
 G_i=G_i'=G,\qquad F=(n+1)G,\qquad
 0\le\dot G\le1,\quad G(1)=\frac n{n+1},\quad
 (n+1)G(s)\ge ns.
\end{equation}
\end{lemma}

\begin{proof}
In the exact reference basis, normalize the $i$th coefficient of the
direction of $T_i'$ to one.  Replacement transversality and Cramer's
rule bound this normalization and all ratios by subpowers.  Include in
the joint limiting law the clipped inverse-size exponents
\[
 r_j=\lim\min\{1,\max\{0,-\log_L|a_j|\}\},\qquad j\ne i,
\]
with value one when a ratio vanishes.  At most one of these $n$ exponents
can be positive.  Otherwise two coefficients are small at a fixed
positive-power scale, placing $T_i'$ in a codimension-two frequency
strip by Lemma~\ref{t:rank}.  Conditional on $X$ and the exact references,
the alternative still has its original law.  Conditional independence
and \eqref{t:strip-law} make this event have probability tending to zero.

On $\{r_j=0\}$ compare profile increments using the $j$th coordinate.
For fixed rational $0<a<t/2$, the finite event $r_{j,L}<a/2$ implies
$|a_j|>L^{-a/2}$.  The approximate basis fixed only by the reference
bins changes this to a lower bound $L^{-a}$ for large $L$.
Given $T_i',Z_j,C,D$, the child position then has at most
$L^{O_n(a)+o(1)}$ possibilities, since that coordinate parametrizes
the tube with inverse slope at most $L^{a+o(1)}$.  This permitted set of
child labels is determined by the displayed $(T_i',Z_j,C,D)$, uniformly
over hidden exact reference directions.
Apply restricted determination before taking limits, and then let
$a\downarrow0$.  On $\{r_j=0\}$ it gives
\[
 G_i(s+t)-G_i(s)-\ell_i'
       \le I(Z_j\mid C,D)
       \le G_j(s+t)-G_j(s)+O_n\left(\sum_k\ell_k\right).
\]
The comparison uses only the indicated bin-defined coordinates; exact
reference directions have not been silently added to the conditioning.

On dyadic exponent partitions with block size $b=b_*$, distribute each
cost uniformly over its block.  Its expected integral is $O_n(b)$.
Dyadic summability makes all these cost densities tend to zero for
almost every path and exponent.  At simultaneous differentiability
points, $\dot G_i\le\dot G_j$ whenever $r_j=0$.  These ratio events
belong to the entire sampled path, independently of the current block.
All endpoint integrals equal $n/(n+1)$, so each such derivative
inequality implies equality of the two profiles.  Each profile thus
agrees with at least $n-1$ others.  Every equivalence class has at least
$n$ members, and two disjoint classes would need $2n>n+1$ members.
There is a single class, proving \eqref{t:common-profile}.
\end{proof}

Use family one and two specified off-axis ratios $a_2,a_3$ of $T_1'$.
Conditional on the exact $X$ and exact references, redraw $T_1'$ from
its original law.  Let $h_j(u)$ be the limiting negative base-$L$
logarithm of the mass of the radius-$L^{-u}$ ball centered at its sampled
$a_j$.  The functions are nonnegative and nondecreasing and satisfy
\begin{equation}\label{t:angular-sum}
 h_2(u)+h_3(u)\ge c_1u\qquad(0<u<1/8,\ u\in\mathbb Q).
\end{equation}
Indeed the joint ball pulls back to a codimension-two strip by
Lemma~\ref{t:rank}, and the original conditional law obeys
\eqref{t:strip-law}.  Joint grid information is at most the sum of
marginal grid informations by the chain rule and limiting conditioning.
Centered-ball and grid exponents agree: if $p_z$ is a sampled grid-cell
mass and $q_z$ the sum over its boundedly many neighbors, then
\[
 \sum_{z:q_z>L^a p_z}p_z
       \le L^{-a}\sum_zq_z\lesssim L^{-a}.
\]
Apply this to joint and marginal grids, with a boundedly finer grid
where necessary.  This is the ``Angular marginal information'' argument
of \citet[Section~3]{PSE2026}; Lemma~\ref{t:rank} supplies its new
higher-dimensional hypothesis.

\subsection{A planar projection test with no frozen extra coordinates}

We record the precise external projection statement.

\begin{lemma}[Robust planar projection]\label{t:projection}
Fix $0<d<1$ and $\rho>0$.  There are $\eps_{\rm pr},\eta_{\rm pr}>0$
such that, for sufficiently small $\Delta$, the following holds.
Let $A\subset B(0,1)\subset\R^2$, and write $\mathcal N_\Delta$ for
covering number at radius $\Delta$.  Suppose
\[
 \Delta^{-2d+\eta_{\rm pr}}\le\mathcal N_\Delta(A)
       \le\Delta^{-2d-\eta_{\rm pr}},\qquad
 \mathcal N_\Delta(A\cap B(z,r))
       \le\Delta^{-\eta_{\rm pr}}r^\rho\mathcal N_\Delta(A)
\]
for $\Delta\le r\le1$.  If a probability $\sigma$ on projection lines
obeys $\sigma(B(V,r))\le\Delta^{-\eta_{\rm pr}}r^\rho$ on the same
range, then some $V\in\supp\sigma$ satisfies
\[
 \mathcal N_\Delta(\pi_VA')\ge\Delta^{-d-\eps_{\rm pr}}
 \quad\text{for every }A'\subset A\text{ with }
 \mathcal N_\Delta(A')\ge\Delta^{\eta_{\rm pr}}\mathcal N_\Delta(A).
\]
The tolerance and gain may be decreased.
\end{lemma}

\begin{proof}
This is \citet[Theorem~1]{He2020} with ambient dimension $2$,
projection dimension $1$, and set exponent $2d$.  Its Grassmannian
nonconcentration condition is angular-ball nonconcentration in this
case.  The theorem gives the conclusion for all large subsets
simultaneously, outside a direction set of probability a positive power
of $\Delta$.  Decrease the tolerance to obtain the stated version.
\end{proof}

\begin{lemma}[Projection block test]\label{t:block-test}
Fix $0<d<1$, $d_0>0$, and $\kappa_0>0$.  A sufficiently small
$\zeta>0$, depending on these parameters and fixed $n$, has the
following property.  Fix rational $[s,s+b]\subset[0,1]$, $b<1/8$,
$b_*\ge b$, and $j\in\{2,3\}$.  Let $D_0$ be the reference direction
bins at precision $L^{-b_*}$, and $D_1$ additionally contain the bin
of $T_1'$.  The following cannot hold on an event of positive limiting
probability:
\begin{enumerate}[label=\textnormal{(\roman*)}]
\item $|G(s+b)-G(s)-db|\le\zeta b$, and
      $G(s+u)-G(s)\ge d_0u-\zeta b$;
\item every cost for $D_0,D_1$, the reference tubes, and $T_1'$, at the
      endpoint and the tested prefixes, is at most $\zeta b$;
\item $h_j(u)\ge\kappa_0u-\zeta b$.
\end{enumerate}
The prefix tests need only use a fixed finite rational grid in $0<u\le b$
containing $b$, of mesh a sufficiently small fixed multiple of $\zeta b$.
All error constants and relative meshes are independent of $s,b,b_*$.
\end{lemma}

\begin{proof}
Set $K=L^b$, and use the same reference-bin basis for $D_0,D_1$.
Let $Y_u=(Z_1,Z_j)$ at prefix $u$.  Lemmas~\ref{t:additivity}
and~\ref{t:common} give, with either bin conditioning,
\begin{equation}\label{t:pair-information}
 I(Y_b\mid C,D_a)=2db+O_n(\zeta b),\qquad
 I(Y_u\mid C,D_a)\ge2d_0u-O_n(\zeta b),\quad a=0,1.
\end{equation}
Let $\widetilde a_j$ be the alternative ratio determined by $D_1$,
and let $\Pi$ be the interval label at resolution $K^{-1}$ of
$z_j-\widetilde a_jz_1$.  Given $T_1',C,D_1$, it has $L^{o(1)}$
possible values: the central line is part of the tube label, the slope
error is $L^{-b_*+o(1)}$, and the relative width is
$L^{s-1+o(1)}\le L^{-b+o(1)}$.  Moreover coordinate one parametrizes
the full child position along this tube to subpower ambiguity.  Thus
\[
 I(Y_b\mid\Pi,C,D_1)
 \ge I(Y_b\mid\Pi,C,D_1,T_1')
 =I(X_{s+b}\mid C,D_1,T_1')
 \ge db-O_n(\zeta b).
\]
Since $\Pi$ is determined by $Y_b,D_1$, the chain rule and
\eqref{t:pair-information} imply
\begin{equation}\label{t:projection-information}
 I(\Pi\mid C,D_1)\le db+O_n(\zeta b).
\end{equation}
Every statement so far conditions on parent and direction bins, not
on the unused coordinates.

Pass the finite test to finite $L$ with slack.  It succeeds with
probability at least a fixed $p_*>0$ along a subsequence.  Fix parent
and reference-bin labels whose conditional success probability
$p\ge p_*$, and write $\mu$ for that conditional law.  No lower bound
on the probability of this individual fiber is needed.  Absorb fixed
slack into $\alpha=C_n\zeta+o(1)$.  If $A_{\rm bin}$ denotes the
alternative bin, success implies the following bounds for the displayed
labels whenever they have a successful realization:
\begin{equation}\label{t:fiber-masses}
 \begin{split}
 K^{-2d-\alpha}\le\mu(Y_b=y)&\le K^{-2d+\alpha},\\
 K^{-2d-\alpha}\le\mu(Y_b=y\mid A_{\rm bin}=a)
       &\le K^{-2d+\alpha},\\
 \mu(Y_u=y_u)&\le K^\alpha L^{-2d_0u},\\
 \mu(\Pi=\pi\mid A_{\rm bin}=a)&\ge K^{-d-\alpha}.
 \end{split}
\end{equation}
Let $A$ be the centers of all endpoint pair cells with a successful
realization in this fixed fiber.  The first line gives
\begin{equation}\label{t:common-planar-set}
 pK^{2d-\alpha}\le|A|\le K^{2d+\alpha}.
\end{equation}
The mass upper bound for a prefix divided by the endpoint mass lower
bound counts the successful descendants in that prefix.  A ball is
covered by boundedly many cells at a comparable coarser tested prefix;
the finite mesh costs $K^{O_n(\zeta)}$.  Hence
\begin{equation}\label{t:planar-nonconcentration}
 \frac{\mathcal N_{K^{-1}}(A\cap B(z,r))}
      {\mathcal N_{K^{-1}}(A)}
 \lesssim p^{-1}K^{O_n(\zeta)+o(1)}r^{d_0},
                 \qquad K^{-1}\le r\le1.
\end{equation}
The direct exponent is $2d_0$; weakening it allows all grid and short
prefix errors to be absorbed.

Retain alternative bins whose conditional success probability is at
least $p/2$.  Their successful mass is at least $p/2$.  For each such
bin, let $A_a\subset A$ be its successful endpoint pair cells.  The
second and fourth lines of \eqref{t:fiber-masses} give
\begin{equation}\label{t:large-subsets}
 |A_a|\ge(p/2)K^{2d-\alpha},\qquad
 \mathcal N_{K^{-1}}(\pi_{\widetilde a_j(a)}A_a)
       \le K^{d+O_n(\zeta)+o(1)}.
\end{equation}
Here initially $\pi_a(z_1,z_j)=z_j-az_1$; normalizing its projection
vector can only decrease image distances.

We now construct one nonconcentrating direction measure without paying
an inverse probability for an individual alternative bin.  Restrict
actual samples to success and the retained bins, normalize by their
mass (at least $p/2$), and push them to the projection line spanned by
$(-\widetilde a_j,1)$.  To bound a slope ball, first condition on the
exact point and exact references.  They determine the fixed parent and
reference-bin fiber, and the alternative law is still the original one
given the point.  If the slope ball meets success, choose a successful
alternative in it.  Its centered-mass test (iii), at a fixed enlarged
radius, bounds the original probability of the whole ball.  Rounding
to a coarser tested radius costs $K^{O_n(\zeta)+o(1)}$.  Restriction
only decreases this probability.  Integrating over exact points and
references, and finally normalizing the entire successful mass, gives
\begin{equation}\label{t:angular-probability}
 \sigma(B(V,r))\lesssim p^{-1}K^{O_n(\zeta)+o(1)}r^{\kappa_0},
                    \qquad K^{-1}\le r\le1.
\end{equation}
Subpower slope bounds, condition numbers, and the sine-angle formula
justify the passage from slopes to projective lines.  The approximate
ratios have error $K^{-1}L^{o(1)}$.  Rescale the coordinate range of
$A$, of size at most $L^{o(1)}$, into a unit ball; covering numbers and
all displayed bounds change by subpowers only.

Choose $0<\rho<\min\{d_0,\kappa_0\}$.  Lemma~\ref{t:projection}
applies to \eqref{t:common-planar-set}, \eqref{t:planar-nonconcentration},
and \eqref{t:angular-probability} when $\zeta$ is sufficiently small.
For every direction in $\supp\sigma$, however, \eqref{t:large-subsets}
provides a sufficiently large subset of the same $A$ with projection
smaller than the guaranteed gain.  This is a contradiction.  These
are exactly the pair-cell, projection-cell, cost, and angular hypotheses
of the ``Projection block test'' in \citet[Section~3]{PSE2026}; the
proof above verifies their dimensional transfer explicitly.
\end{proof}

\subsection{From the block test to a full-to-empty transition}

The remaining profile argument is one-dimensional in the scale
parameter.  We give its hypotheses and mechanism explicitly, so that
the replacement of the endpoint fraction $2/3$ by $n/(n+1)$ is visible.
The angular sum bound \eqref{t:angular-sum} does not select one marginal
that is nonconcentrated at every radius.  The following finite choice
of comparable scales supplies the angular bounds needed by the block
test.

\begin{lemma}[Finite angular alternatives]\label{t:angular-alternatives}
For fixed $d\in(0,1)$ and $d_0>0$, there is a fixed finite list of
positive rational ratios, bounded away from zero, with this property.
For every sufficiently small rational $b_0$, almost every path admits
a length $b'$ from that list times $b_0$, and $j=2$ or $3$, satisfying
the angular hypothesis of Lemma~\ref{t:block-test}, with spare slack,
for all tested block lengths in $[b'/4,b']$.  Each alternative has a
fixed positive angular slope, tolerance, and finite relative mesh,
independent of $b_0$.
\end{lemma}

\begin{proof}
Use the block-test tolerance $\zeta_2$ for angular slope $c_1/4$.
Choose a positive rational $a\ll c_1\zeta_2$, and then a sufficiently
reduced tolerance $\zeta_1$ for slope $a/2$.  Test
$h_2(u)\ge au$ on a fine fixed relative grid from $c\zeta_1b_0$ to
$b_0$.  If every test passes, monotonicity fills gaps and nonnegativity
handles smaller radii with the allowed additive error; choose $j=2$,
$b'=b_0$.  Otherwise a tested $v$ obeys $h_2(v)<av$.  By monotonicity
and \eqref{t:angular-sum},
\[
 h_3(u)\ge c_1u-h_2(u)>c_1u-av\ge(c_1/2)u
          \quad(2av/c_1\le u\le v).
\]
Choose $j=3$, $b'=v$.  The choice $a\ll c_1\zeta_2$ makes the omitted
smaller radii harmless for every block length in $[v/4,v]$.
Reducing cutoffs and meshes leaves spare slack.  The finitely many
$v/b_0$ are fixed and bounded away from zero.  This proves the
``Finite angular alternatives'' lemma of \citet[Section~3]{PSE2026}
from precisely its two monotone-marginal and positive-sum hypotheses.
\end{proof}

\begin{lemma}[Binary derivative and finite alternation]\label{t:alternation}
Almost surely $\dot G\in\{0,1\}$ almost everywhere, and
$E=\{s:\dot G(s)=1\}$ agrees modulo null sets with a finite union of
intervals.  There is an interior full-to-empty transition.  Consequently,
for every fixed $\eta>0$, some rational block $[s,s+b]$ with $0<b<1/8$
satisfies, on an event of positive probability,
\begin{equation}\label{t:jump}
 G(s+b)-G(s)=b/2+O(\eta b),\qquad
 G(s+b/2)-G(s)=b/2+O(\eta b).
\end{equation}
\end{lemma}

\begin{proof}
We use the arguments named ``Binary derivative'', ``Balanced subintervals
inside a hill'', and ``Finite alternation and a full-to-empty jump'' in
\citet[Section~3]{PSE2026}.  Their abstract hypotheses here are:
$G$ is $1$-Lipschitz and nondecreasing; \eqref{t:angular-sum} holds;
Lemma~\ref{t:block-test} holds with fixed positive tolerances for each
$d,d_0,\kappa_0$; and costs on disjoint block lists total $O_n(b_*)$ in
expectation.  All have been proved.  The following details show that
neither the number of coordinates nor the planar endpoint value enters
these arguments.

If interior derivative values occur on a set of positive
path-times-exponent measure, choose $d\in(0,1)$ and $0<d_0<d$ so
that a small neighborhood of $d$ captures positive measure.  At dyadic
$b_0\downarrow0$, form a rational mesh of candidate blocks for every
alternative in Lemma~\ref{t:angular-alternatives}.  Endpoints have spacing
$qb_0$, with fixed rational $q$ below all relative tolerances; lengths
range between fixed positive multiples of $b_0$.  Include each finite
prefix mesh and use direction precision $b_*=b_0$.  Each family has
bounded overlap, so can be colored into finitely many disjoint lists.
If $\mathcal C_I$ is the sum of all costs for candidate $I$, then
\begin{equation}\label{t:mesh-cost}
 \sum_I\mathbb E\mathcal C_I\lesssim_n b_0.
\end{equation}
For $R_{b_0}(\omega,s)=\sum_{I\ni s}\mathcal C_I(\omega)/b_0$,
integration over $s$ gives $\mathbb E\int R_{b_0}\lesssim_n b_0$.
Dyadic summability forces $R_{b_0}\to0$ almost everywhere.
At differentiability points, all nearby increments and tested prefixes
have their linear approximation with error $o(b_0)$.  The angular
alternative and the vanishing costs then give a successful block test.
The rational meshes and test parameters form a countable family.  For
each fixed test, Lemma~\ref{t:block-test} makes its all-success event
null; the preceding positive-measure interior-derivative event would
force one of those null events to have positive probability.  A countable cover
of interior derivative values proves that the derivative is binary.

We also need a deterministic fact about this binary derivative.  If
$a<b$ are respectively density points of $E$ and $E^c$, then for every
sufficiently small $r>0$ there is $[x,y]\subset(a,b)$ with
\begin{equation}\label{t:balanced-hill}
 r\le y-x\le2r,\qquad
 G(y)-G(x)=(y-x)/2,\qquad
 G(x+u)-G(x)\ge u/2\quad(0\le u\le y-x).
\end{equation}
To see this, set $H(s)=G(s)-s/2$.  The density assumptions give a
closed hill strictly inside $(a,b)$, with equal endpoint heights and
larger interior values.  Among levels admitting a closed subinterval
of length at least $r$ with values above the level inside, choose the
highest, by compactness.  No component strictly above this level has
length greater than $r$, since it would admit a higher such interval.
The level set has no gap longer than $r$.  Starting at the left
endpoint, its first point at distance at least $r$ lies at distance
between $r$ and $2r$, and gives \eqref{t:balanced-hill}.  Finitely many
disjoint ordered density pairs give disjoint balanced intervals at
every sufficiently small common $r$.

If a positive-probability set of paths had arbitrarily many disjoint
ordered full/empty density pairs, apply the angular alternatives with
$d=1/2$, $d_0=1/4$.  For a successful angular size $b'\asymp b_0$, use
\eqref{t:balanced-hill} at a size yielding lengths in $[b'/3,2b'/3]$,
and round endpoints to the fixed fine relative mesh.  Any prescribed
number of disjoint hills gives that many spatially and angularly
successful candidates for all sufficiently small $b_0$.  Each must
fail a cost test, by Lemma~\ref{t:block-test}; such a failure costs a
fixed positive multiple of $b_0$.  Equation~\eqref{t:mesh-cost} bounds
the expected number of failures uniformly in $b_0$.  Fatou's lemma
contradicts its tending to infinity on a positive-probability set.
Thus almost every path has finitely many alternations among density
points.  Their ordered types, together with the Lebesgue density
theorem, make $E$ a finite union of intervals modulo null sets.

Finally \eqref{t:common-profile} gives
$|E|=n/(n+1)<1$ and $(n+1)|E\cap[0,s]|\ge ns$.  The latter prevents
an initial empty interval and the former forces a later empty interval.
There is therefore an interior full-to-empty jump.  Choose a small
block with midpoint within $\eta b$ of the jump, round its endpoints
rationally, and keep its halves inside the adjacent full and empty
intervals except for this error.  The countable rational choices cover
almost every path, so one has positive probability.  This gives
\eqref{t:jump} with $b<1/8$.
\end{proof}

\subsection{The growing-degree partitioning contradiction}

\begin{lemma}[No full-to-empty block]\label{t:no-jump}
For sufficiently small fixed $\eta>0$, the positive-probability event
in \eqref{t:jump} is impossible.
\end{lemma}

\begin{proof}
Fix the rational block and put $K=L^b$.  Endpoint and midpoint
increments for $F$ are both $(n+1)b/2+O_n(\eta b)$; the conditional
endpoint increment given $T=T_1$ is $b/2+O(\eta b)$.
Pass to finite scale with slack and fix a parent $C=X_s=c$ having
conditional success probability at least a fixed $p_*>0$.  Write
\[
 p_c(z)=\mathbb P(X_{s+b}=z\mid C=c),\qquad
 q_T(z)=\mathbb P(X_{s+b}=z\mid T,C=c).
\]
Call $(T,z)$ good when it has some successful realization.  Let
$\mathcal Z$ be the good endpoint labels, represented by their centers
in normalized parent coordinates.  Since the bounds depend only on
the displayed labels, goodness by existence preserves
\begin{equation}\label{t:good-masses}
 K^{-(n+1)/2-O_n(\eta)}\le p_c(z)
       \le K^{-(n+1)/2+O_n(\eta)},\qquad
 q_T(z)\le K^{-1/2+O_n(\eta)}\quad\text{for good }(T,z).
\end{equation}
The midpoint mass upper bound divided by the endpoint mass lower
bound gives
\begin{equation}\label{t:good-counts}
 |\mathcal Z|\le K^{(n+1)/2+O_n(\eta)},\qquad
 \#(\mathcal Z\cap B(z,CK^{-1/2}))\le K^{O_n(\eta)}.
\end{equation}
Good incidences have central-line error
$\delta=K^{-1}L^{o(1)}$ in this bounded parent.

Choose small fixed $\eps>0$ and $P=K^{1/2-\eps}$.  Polynomial
partitioning in $\R^{n+1}$ gives a nonzero polynomial of degree $O_n(P)$
whose open complementary cells each contain at most
\begin{equation}\label{t:partition-occupancy}
 O_n(|\mathcal Z|/P^{n+1})=K^{(n+1)\eps+O_n(\eta)}
\end{equation}
points; points on the zero set are assigned to the wall.  This is the
finite-point polynomial partition theorem
\citep[Theorem~4.1]{GuthKatz2015}, whose polynomial ham-sandwich
construction works in every fixed ambient dimension.

Take a wall of thickness a sufficiently large constant times $\delta$.
At midpoint resolution $r=K^{-1/2}$, Wongkew's hypersurface
neighborhood-volume bound \citep[Main Theorem]{Wongkew1993}, in a fixed
ball slightly larger than the parent, gives the number of boxes meeting
this wall as
\begin{equation}\label{t:wall-count}
 O_n\left(\sum_{a=1}^{n+1}P^a r^{a-(n+1)}\right)
       =O_n(PK^{n/2}).
\end{equation}
Indeed $\delta=o(r)$, so thicken to $O(r)$, use the volume estimate,
and divide by $r^{n+1}$.  The equality follows from $Pr=K^{-\eps}$.
The formula includes singular zero sets and all degree powers; the
degree is allowed to grow here.  Equations~\eqref{t:good-masses}
and~\eqref{t:good-counts} bound the discarded good mass by
\begin{equation}\label{t:wall-mass}
 O_n(PK^{n/2})K^{O_n(\eta)}K^{-(n+1)/2+O_n(\eta)}
       =K^{-\eps+O_n(\eta)}=o(1),
\end{equation}
once $\eta\ll_n\eps$.

The wall contains the incidence error, so the nearby central-line point
of each surviving endpoint belongs to the same open cell.  A line not
contained in the zero set crosses at most $O_n(P)$ cells; a line in it
has no off-wall incidence.  Draw $T$ with its original marginal given
$C=c$, and then $z,z'$ independently with law $q_T$.  If $g(T)$ is
the $q_T$ mass of good off-wall endpoints, then
$\mathbb E g(T)\ge p_*/2$.  Cauchy over the cells visited by $T$,
followed by Jensen, gives
\begin{equation}\label{t:pair-lower}
 \mathbb P\{z,z'\text{ good off-wall in the same cell}\mid C=c\}
       \gtrsim_n P^{-1}=K^{-1/2+\eps}.
\end{equation}

For $|z-z'|\lesssim K^{-1/2}$ there are at most $K^{O_n(\eta)}$
possible good second endpoints by \eqref{t:good-counts}; each costs
$K^{-1/2+O_n(\eta)}$ under $q_T$.  Near pairs therefore contribute
at most $K^{-1/2+O_n(\eta)}$.

For a farther pair, a line incident to both endpoints to error $\delta$
has direction in one of two caps of radius $K^{-1/2+o(1)}$, hence in
caps of radius $K^{-1/3}$ for large $L$.  Under the double sampling,
the $(T,z)$ marginal is the original joint law.  Thus
\begin{equation}\label{t:posterior-law}
 \mathbb P(T\in\mathcal A\mid z,C=c)
 =\sum_{x:X_{s+b}(x)=z}\mathbb P(X=x\mid z,C=c)
                       \mathbb P(T\in\mathcal A\mid X=x).
\end{equation}
An angular cap is contained, on bounded frequency sets, in a
codimension-two strip of comparable width.  Use \eqref{t:strip-law}
at exponent $b/3<1/8$ in each term of this mixture.  Its cap probability
is at most $K^{-c_1/3+o(1)}$.
For each first endpoint, the same-cell list has at most
$K^{(n+1)\eps+O_n(\eta)}$ second endpoints by
\eqref{t:partition-occupancy}; this list is fixed before drawing $T$.
On a good second incidence, bound $q_T(z')$ by
\eqref{t:good-masses}, then apply the cap bound to the remaining
indicator and sum the list.  The far-pair contribution is at most
\[
 K^{-1/2-c_1/3+(n+1)\eps+O_n(\eta)+o(1)}.
\]
No angular estimate conditional on the success event has been used.
After division by \eqref{t:pair-lower}, the near and far bounds are
\[
 K^{-\eps+O_n(\eta)},\qquad
 K^{-c_1/3+n\eps+O_n(\eta)+o(1)}.
\]
Choose $\eps\ll_n c_1$, then $\eta\ll_n\eps$.  Both tend to zero,
contradicting \eqref{t:pair-lower}.
\end{proof}

\begin{proof}[Proof of Proposition~\ref{t:transverse}]
A countersequence yields Lemma~\ref{t:tubes}.  Its information profiles
coincide by Lemma~\ref{t:common}; Lemma~\ref{t:alternation} supplies
a full-to-empty block with any prescribed small fixed tolerance.
Lemma~\ref{t:no-jump} forbids it.  Hence no countersequence exists,
which supplies a fixed $c>0$ and a large-scale threshold as asserted.
\end{proof}
